\chapter{Conditional fitness, reproductive contributions and logarithms}
\label{conditioning:chapter}

An individual's breeding value does not determine its realised reproductive
output. Individuals with the same breeding value can differ in environment,
genetic background and unmeasured traits, and reproduction itself can be
stochastic. We therefore begin with an average over the individuals represented
at each breeding value. This chapter derives what that average supplies, why its
logarithm is taken after averaging, and how environmental variation changes it.
The variance calculations in Chapter~\ref{gaussian:chapter} then describe
variation in this specified quantity.

\section{A conditional mean defines the fitness relationship}

Let \(\vect A\) denote breeding value in a specified population and let \(W\)
denote nonnegative reproductive output during one selection episode. Assume
\[
 0<\bar W=\E[W]<\infty.
\]
The conditional expected output is
\begin{equation}
 W_A(\vect a)=\E[W\mid\vect A=\vect a].
 \label{conditioning:target}
\end{equation}
The notation describes an expectation under the current joint distribution of
breeding values and reproductive output. It does not require a deterministic
map from breeding value to offspring count. Conditioning on a continuous
variable is understood through a conditional distribution, rather than through
the probability of the single event \(\vect A=\vect a\).

The \emph{tower property} states that an unconditional average can be obtained
by first averaging conditionally and then averaging those conditional means.
For example,
\[
 \E[W_A(\vect A)]
 =\E\{\E[W\mid\vect A]\}
 =\E[W]=\bar W.
\]
More generally, if \(f(\vect A)\) is a scalar quantity with
\(\E[|f(\vect A)|W]<\infty\), then
\begin{align}
 \E[f(\vect A)W]
 &=\E\{\E[f(\vect A)W\mid\vect A]\}\nonumber\\
 &=\E\{f(\vect A)\E[W\mid\vect A]\}\nonumber\\
 &=\E[f(\vect A)W_A(\vect A)].
 \label{conditioning:tower_weight}
\end{align}
The second line follows because \(f(\vect A)\) is already known when
\(\vect A\) is given. This identity is the reason conditional expected output
can be used to calculate the distribution of parental breeding values
represented by reproductive contributions.

\begin{proposition}[Expected reproductive weighting]
\label{conditioning:selection_weight}
Under the assumptions above, define a probability law \(Q\) by
\[
 Q(C)=\frac{\E[{\boldsymbol 1}_{\{\vect A\in C\}}W]}{\bar W}
\]
for each set \(C\) of breeding values. Its expectation of \(f\) is
\begin{equation}
 \E_Q[f(\vect A)]
 =\frac{\E[f(\vect A)W_A(\vect A)]}{\bar W}.
 \label{conditioning:weighted_expectation}
\end{equation}
If the initial distribution has density \(p_A\), the weighted density is
\begin{equation}
 p_{\rm sel}(\vect a)
 =\frac{W_A(\vect a)p_A(\vect a)}{\bar W}.
 \label{conditioning:selected_density}
\end{equation}
\end{proposition}

\begin{proof}
Nonnegative output makes \(Q(C)\) nonnegative. Taking \(C\) to be the entire
space gives \(Q(C)=\E[W]/\bar W=1\), and countable additivity follows from
additivity of the expectation of nonnegative indicators. Thus \(Q\) is a
probability law. Equation~\eqref{conditioning:tower_weight} gives its
expectation of \(f\). When a density exists, write the numerator as
\[
 \E[f(\vect A)W_A(\vect A)]
 =\int f(\vect a)W_A(\vect a)p_A(\vect a)\,d\vect a.
\]
Dividing by \(\bar W\) identifies the density in
Equation~\eqref{conditioning:selected_density}. Its integral is one because
the numerator integrates to \(\E[W_A(\vect A)]=\bar W\).
\end{proof}

The construction remains valid without a density. For two classes with
frequencies \(p_1,p_2\) and expected outputs \(w_1,w_2\), their reproductive
fractions are \(p_jw_j/(p_1w_1+p_2w_2)\). If the outputs are one and two,
increasing the frequency of the second class from \(1/4\) to \(3/4\) changes
mean output from
\[
 \tfrac34(1)+\tfrac14(2)=\tfrac54
 \quad\hbox{to}\quad
 \tfrac14(1)+\tfrac34(2)=\tfrac74.
\]
Its reproductive fraction consequently changes from
\[
 \frac{(1/4)2}{5/4}=\frac25
 \quad\hbox{to}\quad
 \frac{(3/4)2}{7/4}=\frac67.
\]
The expected output at each breeding value has not changed. Its frequency has.

Taking \(f(\vect A)=\vect A\) component by component also gives a general
selected-parent displacement. If the required first moments are finite,
\begin{align}
 \E_Q[\vect A]-\E[\vect A]
 &=\frac{\E[\vect A W]-\E[\vect A]\E[W]}{\bar W}\nonumber\\
 &=\frac{\Cov(\vect A,W)}{\bar W}
 =\frac{\Cov\{\vect A,W_A(\vect A)\}}{\bar W}.
 \label{conditioning:covariance_displacement}
\end{align}
This is the covariance part of the selection accounting developed by
\citet{Price1970}. It uses no Gaussian approximation. Gaussian selected
moments require the additional assumptions introduced in
Chapter~\ref{selection:chapter}.

\paragraph{Biological interpretation.}
The weights describe how parental breeding values are represented among
expected reproductive contributions. They do not specify an offspring's
breeding value. That requires a transmission model. Nor should a population
expectation be confused with the random ratio
\(\sum_i f(\vect A_i)W_i/\sum_i W_i\) in a finite realised sample: expectation
of that ratio need not equal a ratio of expectations. The population law above
is the target against which such finite-sample quantities can be compared.

\section{What taking the logarithm changes}

Where \(W_A(\vect a)>0\), define
\begin{equation}
 \ell_A(\vect a)=\log W_A(\vect a).
 \label{conditioning:log_target}
\end{equation}
For any two breeding values with positive conditional means,
\begin{align}
 \ell_A(\vect a_2)-\ell_A(\vect a_1)
 &=\log W_A(\vect a_2)-\log W_A(\vect a_1)\nonumber\\
 &=\log\!\left\{\frac{W_A(\vect a_2)}{W_A(\vect a_1)}\right\}.
 \label{conditioning:log_ratio}
\end{align}
Thus equal differences in \(\ell_A\) represent equal ratios of expected output.
A difference of \(0.10\) corresponds to a ratio \(e^{0.10}\), approximately
\(1.105\). The gradient of \(\ell_A\) measures a local proportional change in
expected output per unit change in breeding value:
\[
 \frac{\partial\ell_A}{\partial a_j}
 =\frac{1}{W_A}\frac{\partial W_A}{\partial a_j}.
\]
This is the connection to a log link for expected fitness
\citep{MorrisseyGoudie2022}. It concerns the conditional mean, not a
transformation of each observed count.

Changing the unit of reproductive output multiplies all expected outputs by
one positive constant \(c\). Then
\[
 \log\{cW_A(\vect a)\}=\log c+\ell_A(\vect a).
\]
The gradient, Hessian and variance of \(\ell_A\) are unchanged because the
added term is constant. The reproductive weights are unchanged as well:
\[
 \frac{cW_A(\vect a)p_A(\vect a)}{c\bar W}
 =\frac{W_A(\vect a)p_A(\vect a)}{\bar W}.
\]
Neither result makes variance on the log scale equal to variance on the
original output scale. They describe different functions.

\section{Why averaging and taking logarithms do not commute}

\begin{proposition}[Conditional Jensen inequality]
\label{conditioning:jensen}
Fix a breeding value \(\vect a\). Suppose the conditional output is positive
and has finite positive mean \(m=W_A(\vect a)\). If the conditional expectation
of \(\log W\) is finite, then
\begin{equation}
 \E[\log W\mid\vect A=\vect a]\leq\log W_A(\vect a).
 \label{conditioning:jensen_inequality}
\end{equation}
Equality holds precisely when \(W=m\) with conditional probability one.
\end{proposition}

\begin{proof}
For \(t>0\), put \(h(t)=t-1-\log t\). Its derivative is
\(h'(t)=1-1/t\), so it decreases up to \(t=1\) and increases after \(t=1\).
Its unique minimum is \(h(1)=0\). Therefore
\[
 \log t\leq t-1,
\]
with equality only at \(t=1\). Apply this inequality to \(t=W/m\):
\[
 \log W\leq\log m+\frac{W-m}{m}.
\]
Taking the conditional expectation gives
\[
 \E[\log W\mid\vect A=\vect a]
 \leq\log m+\frac{\E[W\mid\vect A=\vect a]-m}{m}
 =\log m.
\]
The nonnegative difference has expectation zero only if it is zero almost
surely. By the equality condition above, this requires \(W/m=1\) almost surely.
\end{proof}

If zero output has positive conditional probability, \(\log W\) is not a
finite-valued observation. It may be treated as \(-\infty\) in an extended
expectation, but it cannot be used as an ordinary finite response. Meanwhile
\(\log W_A(\vect a)\) remains finite whenever the conditional mean is positive.
For outputs zero and two with equal probability,
\[
 \log\E[W\mid\vect A=\vect a]=\log 1=0,
\]
whereas one of the individual logarithms is not finite. If the conditional
mean itself is zero, nonnegative output is zero almost surely at that breeding
value. Its contribution to the selected-parent law is then zero, and
\(\ell_A\) has no finite value there. These are different boundary cases.

\paragraph{Biological interpretation.}
Variation in counts among individuals with the same breeding value is averaged
before fitness differences across breeding values are described. Excluding zero
counts, adding a constant to every count before logging, or averaging logged
positive counts changes that target. A count likelihood with a log link can
retain zero counts because it models the positive mean while leaving the
observations on their count scale.

\section{When a log observation has the required conditional mean}

The distinction above does allow a useful special case. A log observation can
describe the same shape as \(\ell_A\) if the gap between its conditional mean
and \(\ell_A\) is constant across breeding values.

\begin{proposition}[A constant conditional log gap]
\label{conditioning:log_gap_model}
Assume \(0<W_A(\vect A)<\infty\) almost surely, and suppose
\[
 W=W_A(\vect A)U,
 \qquad U>0,\qquad U\ \hbox{independent of }\vect A,
 \qquad \E[U]=1,\qquad \E[(\log U)^2]<\infty.
\]
Then, with \(Y=\log W\),
\begin{equation}
 Y=\ell_A(\vect A)+\kappa+\varepsilon,
 \qquad
 \kappa=\E[\log U],\qquad
 \E[\varepsilon\mid\vect A]=0,
 \label{conditioning:constant_gap}
\end{equation}
where \(\varepsilon=\log U-\kappa\) has finite variance. Also \(\kappa\leq0\).
\end{proposition}

\begin{proof}
First check that the multiplicative model has the stated conditional mean:
\[
 \E[W\mid\vect A]
 =W_A(\vect A)\E[U\mid\vect A]
 =W_A(\vect A).
\]
Taking a logarithm of this positive observation gives
\[
 Y=\ell_A(\vect A)+\log U.
\]
Independence makes \(\E[\log U\mid\vect A]=\E[\log U]=\kappa\). Subtracting
\(\kappa\) therefore leaves a residual with zero conditional mean. The assumed
second moment makes its variance finite. Finally, Jensen's inequality gives
\(\kappa=\E[\log U]\leq\log\E[U]=0\).
\end{proof}

For an explicit example, take \(Z\sim N(0,1)\) independent of breeding value
and let
\[
 U=\exp(\sigma Z-\sigma^2/2).
\]
Completing the square in the standard normal integral gives
\[
 \E[e^{\sigma Z}]
 =\frac{1}{\sqrt{2\pi}}\int
   e^{\sigma z-z^2/2}\,dz
 =e^{\sigma^2/2}.
\]
Hence \(\E[U]=1\), \(\kappa=-\sigma^2/2\), and
\(\varepsilon=\sigma Z\) has variance \(\sigma^2\).
With \(\sigma^2=4\), the residual variance is four even though the conditional
mean of the log observation differs from \(\ell_A\) by the constant \(-2\).
Under the Gaussian quadratic assumptions, the nested-\(R^2\) result in
Proposition~\ref{gaussian:nested_r2} then applies.

Independence is a sufficient condition, not the only possibility. It is enough
that \(\E[U\mid\vect A]=1\), that
\(\E[\log U\mid\vect A]=\kappa\) is constant, and that the resulting residual
has finite variance. By contrast, if a conditionally mean-one lognormal factor
has variance parameter \(\sigma^2(\vect A)\), then
\[
 \E[\log U\mid\vect A]=-\frac12\sigma^2(\vect A).
\]
When that quantity varies across breeding values, it contributes its own slope
and curvature to \(\E[Y\mid\vect A]\). The required constant-gap model no longer
holds. None of these positive multiplicative models permits taking finite
logarithms of zero offspring counts.

\section{Deriving the environmental average}

We now construct \(W_A\) from a phenotype model. This is an additional
biological assumption, not a property implied by the definition of breeding
value. In one trait, let
\[
 Z=A+E,\qquad E\sim N(0,V_E),\qquad E\ \hbox{independent of }A.
\]
Here \(E\) is an environmental deviation and \(V_E>0\) its variance. Suppose
expected output conditional on the phenotype is
\begin{equation}
 f_Z(z)=W_{\max}
 \exp\!\left\{-\frac{(z-\theta)^2}{2\omega^2}\right\},
 \qquad W_{\max}>0,\quad\omega^2>0.
 \label{conditioning:phenotype_fitness}
\end{equation}
Assume that \(\E[W\mid A,E]=f_Z(A+E)\). This makes phenotype sufficient for
the conditional reproductive mean in this example. The tower property gives
\[
 W_A(a)=\E[f_Z(a+E)\mid A=a].
\]
This is the environmental averaging used in the quantitative-genetic treatment
of selection \citep{Lande1979}.

\begin{proposition}[Gaussian environmental averaging]
\label{conditioning:environmental_average}
Under these assumptions,
\begin{equation}
 W_A(a)=W_{\max}
 \sqrt{\frac{\omega^2}{\omega^2+V_E}}\,
 \exp\!\left\{-\frac{(a-\theta)^2}{2(\omega^2+V_E)}\right\}.
 \label{conditioning:environment_formula}
\end{equation}
At a population mean \(a_0\), its log slope and curvature are
\begin{equation}
 b=\frac{\theta-a_0}{\omega^2+V_E},
 \qquad H=-\frac{1}{\omega^2+V_E}.
 \label{conditioning:environment_derivatives}
\end{equation}
\end{proposition}

\begin{proof}
Put \(\delta=a-\theta\). Substituting the normal density gives
\begin{align}
 W_A(a)
 &=\frac{W_{\max}}{\sqrt{2\pi V_E}}
   \int_{-\infty}^{\infty}
   \exp\!\left\{-\frac12
   \left(\frac{e^2}{V_E}+\frac{(e+\delta)^2}{\omega^2}\right)\right\}\,de.
 \label{conditioning:environment_integral}
\end{align}
To integrate, first expand the expression in parentheses:
\[
 \frac{e^2}{V_E}+\frac{(e+\delta)^2}{\omega^2}
 =\frac{\omega^2+V_E}{V_E\omega^2}e^2
   +\frac{2\delta}{\omega^2}e+\frac{\delta^2}{\omega^2}.
\]
Let
\[
 C=\frac{\omega^2+V_E}{V_E\omega^2},
 \qquad d=\frac{V_E\delta}{\omega^2+V_E}.
\]
The first two terms are \(C(e+d)^2-Cd^2\). The remaining constant is
\begin{align*}
 \frac{\delta^2}{\omega^2}-Cd^2
 &=\frac{\delta^2}{\omega^2}
   -\frac{V_E\delta^2}{\omega^2(\omega^2+V_E)}\\
 &=\frac{\delta^2}{\omega^2+V_E}.
\end{align*}
Therefore
\[
 \frac{e^2}{V_E}+\frac{(e+\delta)^2}{\omega^2}
 =C(e+d)^2+\frac{\delta^2}{\omega^2+V_E}.
\]
The part independent of \(e\) can be taken outside the integral. In the
remaining integral, change variable to \(u=\sqrt C(e+d)\), so
\(de=du/\sqrt C\):
\[
 \int_{-\infty}^{\infty}e^{-C(e+d)^2/2}\,de
 =\frac{1}{\sqrt C}\int_{-\infty}^{\infty}e^{-u^2/2}\,du
 =\sqrt{\frac{2\pi}{C}}.
\]
Multiplying by the original normalising factor gives
\[
 \frac{1}{\sqrt{2\pi V_E}}\sqrt{\frac{2\pi}{C}}
 =\frac{1}{\sqrt{V_EC}}
 =\sqrt{\frac{\omega^2}{\omega^2+V_E}},
\]
which proves Equation~\eqref{conditioning:environment_formula}. Its logarithm
is a constant minus \((a-\theta)^2/\{2(\omega^2+V_E)\}\). Its first derivative
is \(- (a-\theta)/(\omega^2+V_E)\); its second derivative is
\(-1/(\omega^2+V_E)\). Evaluating at \(a_0\) proves
Equation~\eqref{conditioning:environment_derivatives}.
\end{proof}

The expression has a continuous zero-environmental-variance limit:
as \(V_E\) tends to zero, \(W_A(a)\) tends to \(f_Z(a)\).
For the recurring example \(V_E=\omega^2=1\), \(\theta=4\), and \(a_0=0\),
\[
 W_A(a)=\frac{W_{\max}}{\sqrt2}e^{-(a-4)^2/4},
 \qquad b=2,\qquad H=-\frac12.
\]
The phenotype-space curvature is \(-1\); the breeding-value-space curvature
is \(-1/2\). Their difference follows from averaging, before any population
variance calculation is made.

\subsection{The Jensen gap in this specific model}

This example also allows the two orders of operation to be compared
algebraically. Directly averaging the log phenotype-fitness function gives
\[
 \E[\log f_Z(a+E)]
 =\log W_{\max}-\frac{(a-\theta)^2+V_E}{2\omega^2},
\]
because \(\E[E]=0\) and \(\E[E^2]=V_E\). In contrast,
\[
 \log\E[f_Z(a+E)]
 =\log W_{\max}
  -\frac12\log\!\left(1+\frac{V_E}{\omega^2}\right)
  -\frac{(a-\theta)^2}{2(\omega^2+V_E)}.
\]
Subtracting the first expression from the second and writing
\(t=V_E/\omega^2\) gives
\begin{equation}
 \frac12\{t-\log(1+t)\}
 +\frac{(a-\theta)^2V_E}{2\omega^2(\omega^2+V_E)}.
 \label{conditioning:environment_jensen_gap}
\end{equation}
Both terms are nonnegative. For the first, \(\log(1+t)\leq t\) follows from
the same logarithm inequality used above. Unless \(V_E=0\), the gap generally
depends on breeding value through \((a-\theta)^2\). Averaging logged
phenotype-fitness values therefore does not even preserve the breeding-value
curvature in this model.

\paragraph{Biological interpretation.}
Environmental variation exposes individuals with a given breeding value to a
distribution of phenotypes. Their mean output is consequently a different
function from the phenotype-fitness relationship. The variance \(V_E\) in that
average and the breeding-value variance \(G\) used to weight the population
have different roles. Changing \(G\) on a fixed \(W_A\) is the controlled
population comparison analysed later; changing \(V_E\) generally changes
\(W_A\) itself. The width parameters and variances have trait-squared units,
the slope \(b\) has inverse-trait units, and \(H\) has inverse-trait-squared
units.

\section{Connection to the reproducible calculations}

The main and guide figure generators evaluate this environmental model, while
\codepath{code/reproduce_result_checks.py} checks its integral
independently in the group \codepath{conditional_mean_and_log}. That group also
checks the two-class reproductive fractions. Its
\codepath{noisy_regression_ratio} group illustrates the constant-gap observation
model. These are numerical checks of specified examples. The identities above
follow from the stated probability model and algebra; the checks do not
establish them for an unrestricted biological population.
